\begin{table*}[t]
  \centering\small
  \begin{tabularx}{\textwidth}{@{}l>{\raggedright\arraybackslash}Xrrrrrr@{}}
    \toprule
    & & & & \multicolumn{2}{c}{Instruction Count} & \multicolumn{2}{c}{Constraint Satisfaction} \\
    \cmidrule(lr){5-6}\cmidrule(lr){7-8}
    Role & Arm & $T$ & $N$ & $|D_T|$ & Excess Size (\pct) & Rate (\pct, $t{=}T$) & Rate (\pct, $t{=}0$) \\
    \midrule
    \textsc{control} & no prompt comments & 15 & 552 & 3.5 & $+60.4 \pm 18.4$ & $39.0 \pm 2.8$ & $23.2 \pm 2.4$ \\
    \textsc{placebo} & comment-shaped noise & 15 & 552 & 3.3 & $+53.2 \pm 14.6$ & $40.3 \pm 2.7$ & $22.0 \pm 2.4$ \\
    \textsc{treatment} & informative comments & 15 & 552 & 2.1 & {\boldmath$-5.8 \pm 6.6$} & $38.2 \pm 2.6$ & $21.5 \pm 2.4$ \\
    \midrule
    \textsc{control} & no prompt comments & 51 & 184 & 6.6 & $+211.3 \pm 105.3$ & $44.0 \pm 4.9$ & $23.9 \pm 3.9$ \\
    \textsc{placebo} & comment-shaped noise & 51 & 184 & 5.3 & $+147.9 \pm 79.9$ & $42.6 \pm 4.9$ & $22.3 \pm 4.2$ \\
    \textsc{treatment} & informative comments & 51 & 184 & 2.2 & {\boldmath$+1.4 \pm 22.1$} & $44.0 \pm 4.7$ & $21.7 \pm 3.9$ \\
    \bottomrule
  \end{tabularx}
  \caption{
    Informative comments cut excess size at parity satisfaction, more so at $T{=}51$. Excess is
    $|D_T|/|D_\star|{-}1$ at $|D_\star|{=}2.2$; rate is the share of a round's scored
    constraints passed, $t{=}T$ the last-3-round mean. \textbf{Bold} = best in its block, CIs
    disjoint; $\pm$ is half a 95\pct{} bootstrap CI, asymmetric. 3 seeds at $T{=}15$, 1 at $T{=}51$.
  }
  \label{tab:arms}
\end{table*}


\section{Comments Halt the Ratchet and Buy Back Instruction-Following}
\label{sec:lab}

Prompt comments can permanently halt the unbounded growth of instructions in agentic prompts and, in doing
so, improve instruction-following through smaller, more robust prompts.

To evaluate prompt effectiveness, we first propose a new evaluation for agentic prompts in task-stationary
settings (\autoref{sec:testbed}). We use this evaluation to demonstrate that prompt comments enable
prompts to settle near their theoretically optimal minimum cover (\autoref{sec:protocol}). Finally, we show
that prompt comments can buy back instruction-following lost to extraneous, noisy instructions in
real-world prompts (\autoref{sec:wild}).

\subsection{Inverting IFEval makes the minimum cover observable}
\label{sec:testbed}

In practice, maintainers in agentic coding contexts must learn an effective prompt to satisfy a set of
unobservable constraints. In order to maximize correctness, this prompt must simultaneously cover all
constraints while, given well-known issues with instruction-following decay
\citep{jiang2024followbench, lior2025wildifeval}, being minimal.

Evaluating a prompt therefore requires measuring both its excess size and its correctness, but calculating
excess size requires minimum cover, which is incomputable in general (\autoref{sec:theory}). Standard
instruction-following suites offer a path.

By inverting IFEval \citep{zhou2023ifeval}, we can attack the problem. For each benchmark item $(D_j, C_j)$
--- instructions $D_j$, verifiers $C_j$ --- we apply the following transform:

\begin{enumerate}[nosep,itemsep=2pt,parsep=2pt,topsep=2pt,leftmargin=*]
  \item \emph{Hide $D_j$.} The item's stated instructions become the reference minimum cover
    $D_\star$, so $|D_\star|$ is known a priori and excess size is measurable.
  \item \emph{Keep $C_j$.} Its verifiers become the world's hidden constraint set, run by the
    harness alone and never named to the maintainer.
  \item \emph{Generate brief $o_j$.} A stronger model lossily summarizes $D_j$ as a
    general task objective (``draft a product announcement'').
\end{enumerate}

A fresh maintainer then reconstructs $D_j$ from $o_j$ over $T$ steps, seeing only censored, noisy feedback
from $C_j$. Its arm decides whether it may annotate an instruction $d$ with a comment $r_d$ carrying that
instruction's latent reasoning. The next maintainer reads $D_t$ and every $r_d$; the executor reads only $D_t$.

An instruction tells the executor what to do; a comment tells the next maintainer why. Enforcing that split,
the harness strips comments before the prompt reaches the executor and rejects any instruction citing a past
failure, leaving the comment as the only surviving channel for latent reasoning
(\apxref{sec:appendix:testbed}).

\begin{figure}[t]
  \centering
  \definecolor{synComment}{RGB}{0,128,0}
  \scriptsize
  \begin{tabular}{@{}p{0.97\columnwidth}@{}}
    \toprule
    \ttfamily\raggedright Do not truncate or cut off your response; always complete every sentence and thought you begin. \newline
    \textcolor{synComment}{\# r1: response was truncated mid-sentence ("A body that you often s") suggesting response generation stopped prematurely; may indicate token limit, instruction conflict, or assistant aborting output; this directive ensures responses are complete} \tabularnewline
    \midrule
    \ttfamily\raggedright Use the word 'wilderness' exactly 5 times in your response. \newline
    \textcolor{synComment}{\# r3: task 0 failed with 'certain wording is not used the right number of times' — the response uses 'wilderness' only 2 times but requirement expects exactly 5 occurrences; d12 (3 bold sections) was passing, so keeping that approach} \tabularnewline
    \midrule
    \ttfamily\raggedright Divide the response body into exactly 4 paragraphs: (1) introduction to the water cycle, (2) evaporation and condensation, (3) precipitation and collection, (4) conclusion. Separate each paragraph with exactly three blank lines. \newline
    \textcolor{synComment}{\# r3: "response is not split into the right number of paragraphs" recurring across r2-r3; previous directive d18 specified two blank lines between paragraphs but failed; increasing to three blank lines while maintaining strict 4-paragraph structure} \tabularnewline
    \bottomrule
  \end{tabular}
  \caption{
    Each comment names the failure behind its instruction, a hypothesis, and its outcome. Comments
    are green (\texttt{\#}) and never reach the executor. 3 of 8,541 additions, one per
    constraint family.
  }
  \label{fig:comment_examples}
\end{figure}


\subsection{Comments encoding latent reasoning settle the prompt at its cover}
\label{sec:protocol}

Prompt comments encoding an instruction's latent reasoning settle the prompt at its minimum cover
(\autoref{fig:hero:lab}). Its latent reasoning summarizes the failure behind its instruction, a hypothesis,
and how it has fared (\autoref{fig:comment_examples}).

Concretely, across \LabUnits{} maintenance histories, maintainers encoding latent reasoning as prompt
comments learn prompts at \LabPracticalExcessPct{} excess size against \LabControlExcessPct{} for
maintainers without them (\LabPracticalDeltaExcessPp{}), at parity constraint satisfaction
(\autoref{tab:arms}). Arms differ only in the handoff: the
prompt instructions along with its comments (or not), containing a maintainer's summarized latent reasoning.

As (agentic) maintainers scale, they ratchet harder and prompt comments help them more. Varying the
maintainer across \CapacityTiers{} tiers at $T{=}\LabT$, the uncommented arm's excess rises from
\CapacityControlRange{}; at the top tier, prompt comments enable strictly Pareto gains in both constraint
satisfaction and excess size relative to control (\autoref{tab:capacity}).

Finally, ablations at $T{=}\LabT$ show that a comment must carry outcomes (\autoref{tab:ablations}).
Comment-shaped noise lands within the noise of the no-comment arm, and a narrative of attempts without
their outcomes is our worst arm at \LadderFreeformExcessPct{}, handing successors an unvalidated premise to
extend. Within the schema the two fields recording what happened carry the reduction: dropping the recurrence
count alone costs \ClauseDropLossPct{} of it.

\subsection{In real prompts, noisy instructions cost correctness and comments buy it back}
\label{sec:wild}

Extraneous, noisy instructions degrade compliance with the true, correct instructions, and comments
recover most of that loss.

Inverting WildIFEval \citep{lior2025wildifeval} the same way carries the test to real prompts, and to
instruction-following rather than count. We convert every benchmark item $(D_j, C_j)$ to a world as in
\autoref{sec:testbed}, but instead of making the maintainer learn all $|D_j|$ instructions we seed it with
$|D_j|-K$ true instructions and $G$ noisy ones drawn uniformly from other items' sets $D_{i \neq j}$.
WildIFEval's constraints are human-written prose and ship with no code verifiers, so we score them with an arm-blind
LLM judge that sees one constraint and one response and nothing else: no prompt, no arm label, no history.
Across \WildWorlds{} worlds, for $K{=}1$ and $G{=}\WildDistractors$, those noisy instructions cost
\WildInterferencePp{} of correctness on the true instructions already in the prompt, \WildCleanSatPct{} at
$G{=}\WildDistractorsLow{}$ against \WildNoisySatPct{} here (95\% CI: \WildInterferenceCI{}); full
details are in \apxref{sec:appendix:wildgrid}.

Comments lift satisfaction from \WildControlSatPct{} to \WildCommentSatPct{} over \WildT{} maintenance
rounds (\WildDeltaSatPp{}, 95\% CI: \WildDeltaSatCI{}), a \WildDeltaSatPct{} relative gain, against an
uncommented maintainer given the identical prompt. Ablations
show again that the latent reasoning encoded in the comment drives the behavior: comment-shaped noise lands
\WildPlaceboDeltaSatPp{} from the uncommented arm (95\% CI: \WildPlaceboDeltaSatCI{}, covering zero).
Every contrast in this subsection is a rate under one judge, so we re-scored all \JudgeVerdicts{} verdicts
under a second: the criteria reproduce, and the effect that judge measures differs from the one quoted here
by \JudgeDidPp{} (\JudgeAltDeltaSatPp{} against \WildDeltaSatPp{}; 95\% CI: \JudgeDidCI{}), which does
not exclude zero
(\apxref{sec:appendix:judges}).

Under assignment, then, we see that prompt comments encoding latent reasoning both reduce instruction count
to its optimal minimum cover and buy back instruction-following in real-world prompts. Next, we contextualize
and contrast our approach with previous work.
